\chapter*{General Introduction}
\addcontentsline{toc}{chapter}{General Introduction}

The digital transformation of retail banking has pushed the volume and velocity of financial transactions well beyond what manual oversight can monitor. In a network of 150 branches processing thousands of teller operations daily, the gap between transaction speed and supervisory capacity creates exposure to fraud, money laundering, operational errors, and behavioral anomalies that may go undetected until a periodic compliance review surfaces them --- often weeks or months after the fact.

Amen Bank's Direction Centrale du Syst\`eme d'Information already operates an AI-based anti-money laundering system oriented toward bank-wide compliance reporting. This internship addresses a complementary need: \textbf{real-time decision support for branch supervisors at the point of transaction}. The objective is to design and build a platform that ingests teller operations as they occur, profiles both the client and the employee involved, detects anomalies using deep learning models, and surfaces explainable alerts that a non-technical supervisor can understand and act upon immediately.

The platform covers four teller-window operations --- \emph{retrait}, \emph{versement}, \emph{virement}, and \emph{remise de ch\`eques} --- and monitors two independent actors per transaction: the client initiating it and the employee processing it. An Autoencoder detects point anomalies in individual transactions while an LSTM network identifies sequential deviations in client behavior over time. Every alert carries a human-readable explanation generated through SHAP, bridging the gap between model output and supervisory action.

Beyond the detection models themselves, the project invests in the infrastructure that makes a machine learning system operational and maintainable: a streaming pipeline built on RedPanda, behavioral profiles maintained through a hot/cold data layer, experiment tracking with MLflow, three independent retraining triggers, structured logging, and Prometheus-based observability. The system runs as 20 Docker containers orchestrated through Docker Compose, validated end-to-end on a synthetic dataset of 243{,}000 transactions generated from five client archetypes calibrated against the bank's published data.

This rapport is organized as follows. Chapter~\ref{ch:context} presents the institutional and regulatory context and reviews the anomaly detection techniques that inform the design. Chapter~\ref{ch:architecture} describes the architecture, the data schemas, and the key design decisions. Chapter~\ref{ch:synthetic} covers synthetic data generation. Chapter~\ref{ch:training} details model training and evaluation. Chapter~\ref{ch:streaming} presents the streaming pipeline and its integration with the persistence and batch layers. Chapter~\ref{ch:mlops} covers the MLOps layer. Chapter~\ref{ch:dashboard} describes the supervisor dashboard and the end-to-end results. Chapter~\ref{ch:roadmap} documents the enhancements required for a production banking deployment.
